\documentclass[a4paper,10pt]{article}
\usepackage[francais]{babel}
\usepackage{url}

\selectlanguage{francais}
\input{letter-ensiie.sty}

\selectlanguage{francais}
\enCouleur

\begin{document}


\name{St\'ephane \textsc{Genaud}}
\fonction{Directeur de l'antenne de Strasbourg de l'ENSIIE\\Professeur des Universités, ENSIIE}
\telephoneillkirch{+33 (0)3688 54542}
\email{genaud@ensiie.fr}

\begin{letter}[Avis concernant la candidature de Stéphanie Alix à la bourse Eole]{à}%
{Réseau Franco-Néerlandais}

\vspace{-1cm}
Madame, Monsieur,

En tant que directeur du campus de l'ENSIIE à Strasbourg, je souhaite par cette
lettre vous recommander la candidature de Stéphane Alix pour une bourse Eole au
titre d'un stage que Stéphanie souhaite effectuer aux Pays-Bas.

Ce stage  fait partie  du cursus  universitaire de  la formation  d'ingénieur de
l'ENSIIE.  L'ENSIIE est  une école d'ingenieur en informatique formant  en 3 ans
des étudiants recrutés sur un concours  très sélectif après deux années d'études
en  classes  préparatoires aux  mathématiques  et  à  la physique.   Le  diplôme
d'ingénieur  francais  est  aussi  un  grade  Master.  La  commission  du  titre
d'ingénieur en France  préconise que les étudiants passent au  moins deux mois à
l'étranger pendant leurs études.  Nous  soutenons donc fortement les initiatives
comme celles de Stéphanie.
 
Connaissant Stéphanie  Alix depuis son  entrée à  l'école en septembre  2015, et
l'ayant également eu dans mes cours ces deux dernières années, je peux témoigner
de l'extrême  sérieux avec  lequel Stéphanie  mène ses  activités ainsi  que son
travail. Stéphanie  est une très bonne  étudiante, et je suis  convaincu que son
projet de  réaliser son stage  dans l'entreprise \textit{IceMobile}  à Amsterdam
est tout à fait pertinent étant donné ses compétences. L'objet du travail est le
développement d'agents conversationnels (ou \emph{chatbot}) destinés à dialoguer
avec  l'utilisateur dans  le domaine  alimentaire,  agent qui  sera intégré  aux
applications  mobiles ou  web de  l'entreprise. Le  développement d'applications
mobiles est sans  doute l'un des domaines les plus  dynamiques de l'informatique
aujourd'hui,  et permet  à de  nombreuses startups  de faire  émerger des  idées
nouvelles. Ce projet  s'incrit donc dans un  contexte économique potentiellement
prometteur.  Enfin, la  lettre de motivation que Stéphanie m'a  montré me semble
tout à fait  révélatrice de son désir  d'aller à la rencontre de  la culture des
Pays-Bas.


Pour ces raisons, je recommande sans réserve la candidature de Stéphanie Alix
pour la bourse Eole qui pourrait l'aider à concrétiser son stage à Amsterdam.

Veuillez agréer, Madame, Monsieur, l'expression de mes sentiments distingués. 
\end{letter}
\begin{flushright}
\includegraphics[width=.2\textwidth]{signgenaud.jpg}
\end{flushright}
\end{document}
